% Part: incompleteness
% Chapter: incompleteness-provability
% Section: lob-thm

\documentclass[../../../include/open-logic-section]{subfiles}

\begin{document}

\olfileid{inc}{inp}{lob}

\olsection{مبرهنة لوب}

جملة غودل لنظرية $\Th{T}$ نقطة ثابتة للمحمول $\lnot
\OProv[\Th{T}](y)$؛ فهي !!a{sentence} $!G$ تحقق
\[
\Th{T} \Proves \lnot \OProv[\Th{T}](\gn{!G}) \liff !G.
\]
ولا تكون !!{derivable} إذا اتسقت $\Th{T}$.
ذلك أنه لو كان $\Th{T} \Proves !G$ للزم أمران:
(أ) $\Th{T} \Proves \OProv[\Th{T}](\gn{!G})$ من شرط
!!{derivability} الأول؛ و(ب) من $\Th{T} \Proves !G$
ومن $\Th{T} \Proves \lnot \OProv[\Th{T}](\gn{!G})
\liff !G$ يلزم $\Th{T} \Proves \lnot \OProv[\Th{T}](\gn{!G})$،
فتقع $\Th{T}$ في التناقض. فماذا يكون أمر نقطة ثابتة للمحمول
$\OProv[\Th{T}](y)$ نفسه؟ أي !!a{sentence} $!H$ تحقق
\[
\Th{T} \Proves \OProv[\Th{T}](\gn{!H}) \liff !H.
\]
لو كانت !!{derivable}، لأعطى الشرط الأول
$\Th{T} \Proves \OProv[\Th{T}](\gn{!H})$.
وهذه هي النتيجة التي تعطيها قاعدة الفصل من التكافؤ المذكور أيضًا؛
فلا ينتج عدم اتساق $\Th{T}$ بالحجة المستعملة لجملة غودل.
غير أن زوال هذا التناقض لا يثبت أن $\Th{T}$
\emph{تقدر فعلًا} على !!{derive} الصيغة $!H$.

وللجواب نعمّم السؤال قليلًا. أحد اتجاهي التكافؤ السابق،
$\OProv[\Th{T}](\gn{!H}) \lif !H$، حالة من مخطط يسمى
\emph{مبدأ الانعكاس}: $\OProv[\Th{T}](\gn{!A}) \lif !A$.
وسبب التسمية أنه يصوّر النظرية $\Th{T}$ كأنها تتأمل ما يمكنها
!!{derive}، فتقول: «إذا أمكن لـ$\Th{T}$ !!{derive}
الصيغة $!A$، كانت $!A$ صادقة»، أيًّا كانت $!A$.
ولا يصدق هذا إلا للنظريات السليمة؛ فلا يُتوقع من النظرية عمومًا
!!{derive} كل حالة منه. فأي الحالات يمكن لنظرية قوية بالقدر
الكافي ومستوفية لشروط !!{derivability} أن !!{derive}؟
تستطيع ذلك حتمًا حين تكون $!A$ نفسها !!{derivable}؛ وليس وراء
هذه الحالات شيء، كما تقرر المبرهنة الآتية.

\begin{thm}\ollabel{thm:lob}
إذا كانت $!A$ !!a{sentence}، وكانت $\Th{T}$ نظرية
!!a{axiomatizable} تمتد $\Th{Q}$، وكانت
$\OProv[\Th{T}](y)$ صيغة تحقق الشروط P1--P3 الواردة في
\olref[2in]{sec}، فمتى كان النظام $\Th{T}$ !!{derive}s
$\OProv[\Th{T}](\gn{!A}) \lif !A$، لزم أن $\Th{T}$
!!{derive}s الصيغة $!A$ نفسها.
\end{thm}

أي إن $\Th{T} \Proves/ !A$ يستلزم
$\Th{T} \Proves/ \OProv[\Th{T}](\gn{!A}) \lif !A$.
وهذا ما يعرف بمبرهنة لوب.

\begin{explain}
يقرب فكرة البرهان استدلال طريف على وجود سانتا كلوز؛ ومن شاء
نتيجة أخرى وضعها مكان هذه النتيجة. والاستدلال كما يأتي:
\begin{enumerate}
\item نأخذ $X$ الجملة «إذا كانت $X$ صادقة، فإن سانتا كلوز موجود».
\item لنفرض صدق $X$.
\item إذن يصدق مضمونها: إذا صدقت $X$، وجد سانتا كلوز.
\item وبما أننا فرضنا $X$ صادقة، ينتج وجود سانتا كلوز بتطبيق
  قاعدة الفصل على الخطوتين (2) و(3).
\item فاشتققنا (4)، وهي وجود سانتا كلوز، من (2)، وهي صدق $X$.
  وبالبرهان الشرطي حصلنا على «إذا صدقت $X$، وجد سانتا كلوز».
\item وهذا هو مضمون $X$؛ فقد ثبت صدق $X$.
\item نعود إلى الاستدلال (2)--(4)، فينتج وجود سانتا كلوز.
\end{enumerate}
وإذا صغنا هذه الفكرة صوريًا، فجعلنا «!!{derivable}» مكان
«صادقة»، والصيغة $!A$ مكان «سانتا كلوز موجود»، حصلنا على برهان
لوب. ومفتاحه تطبيق لمّة النقطة الثابتة على الـ!!{formula}
$\OProv[\Th{T}](y) \lif !A$؛ فنقطتها الثابتة تقوم مقام
الجملة $X$ في هذا العرض التقريبي.
\end{explain}

\begin{proof}[برهان \olref{thm:lob}]
لتكن $!A$ !!a{sentence}، ولنفرض أن $\Th{T}$ !!{derive}s
$\OProv[\Th{T}](\gn{!A}) \lif !A$. نعرّف $!B(y)$ بأنها
الـ!!{formula} $\OProv[\Th{T}](y) \lif !A$،
ثم تعطي لمّة النقطة الثابتة !!a{sentence} $!D$ بحيث
$\Th{T}$ !!{derive}s $!D \liff !B(\gn{!D})$.
فنحصل، خطوة بعد خطوة، على ما يأتي، وكل منه !!{derivable}
في $\Th{T}$:
\begin{align}
  & !D \liff (\OProv[\Th{T}](\gn{!D}) \lif !A) \ollabel{L-1}\\
  & \qquad \text{إن~$!D$ نقطة ثابتة لـ$!B(y)$}\notag \\
  & !D \lif (\OProv[\Th{T}](\gn{!D}) \lif !A) \ollabel{L-2}\\
  & \qquad\text{من \olref{L-1}}\notag\\
  & \OProv[\Th{T}](\gn{!D \lif (\OProv[\Th{T}](\gn{!D}) \lif !A)}) \ollabel{L-3}\\
  & \qquad \text{من \olref{L-2} بحسب الشرط P1}\notag \\
  & \OProv[\Th{T}](\gn{!D}) \lif \OProv[\Th{T}](\gn{\OProv[\Th{T}](\gn{!D}) \lif !A})
  \ollabel{L-4}\\
  &\qquad \text{من \olref{L-3} باستعمال الشرط P2}\notag \\
  & \OProv[\Th{T}](\gn{!D}) \lif (\OProv[\Th{T}](\gn{\OProv[\Th{T}](\gn{!D})}) \lif \OProv[\Th{T}](\gn{!A})) \ollabel{L-5}\\
  &\qquad \text{من \olref{L-4} باستعمال P2 مرة أخرى} \notag\\
& \OProv[\Th{T}](\gn{!D}) \lif \OProv[\Th{T}](\gn{\OProv[\Th{T}](\gn{!D})}) \ollabel{L-6}\\
  & \qquad\text{بحسب شرط~!!{derivability} P3} \notag\\
  & \OProv[\Th{T}](\gn{!D}) \lif \OProv[\Th{T}](\gn{!A}) \ollabel{L-7} \\
  &\qquad\text{من \olref{L-5} و\olref{L-6}}\notag\\
  & \OProv[\Th{T}](\gn{!A}) \lif !A \ollabel{L-8}\\
  &\qquad\text{بحسب فرض المبرهنة} \notag\\
  & \OProv[\Th{T}](\gn{!D}) \lif !A \ollabel{L-9}\\
  &\qquad\text{من \olref{L-7} و\olref{L-8}}\notag\\
  & (\OProv[\Th{T}](\gn{!D}) \lif !A) \lif !D \ollabel{L-10}\\
  & \qquad \text{من \olref{L-1}}\notag \\
  & !D \ollabel{L-11}\\
  & \qquad\text{من \olref{L-9} و\olref{L-10}}\notag \\
  & \OProv[\Th{T}](\gn{!D}) \ollabel{L-12}\\
  & \qquad\text{من \olref{L-11} بحسب الشرط~P1}\notag \\
  & !A \qquad\qquad\text{من \olref{L-8} و\olref{L-12}}\notag
\end{align}
\end{proof}

وبمبرهنة لوب يقصر برهان عدم الاكتمال الثاني لكل نظرية لها
محمول !!a{derivability} يستوفي P1--P3:
إذا كان $\Th{T} \Proves \OProv[\Th{T}](\gn{\lfalse}) \lif \lfalse$
لزم $\Th{T} \Proves \lfalse$.
لكن اتساق $\Th{T}$ يقتضي $\Th{T} \Proves/ \lfalse$؛
فإذن $\Th{T} \Proves/ \OProv[\Th{T}](\gn{\lfalse}) \lif \lfalse$،
أي $\Th{T} \Proves/ \OCon[\Th{T}]$.
وتجيب المبرهنة عن سؤالنا الأول أيضًا: $!H$، النقطة الثابتة
لـ$\OProv[\Th{T}](x)$، !!{derivable}. فمن
\begin{align*}
  \Th{T} & \Proves \OProv[\Th{T}](\gn{!H}) \liff !H\\
  \intertext{وعلى وجه الخصوص}
    \Th{T} & \Proves \OProv[\Th{T}](\gn{!H}) \lif !H
\end{align*}
ينتج $\Th{T} \Proves !H$ بمبرهنة لوب.

% Going in the other direction, for homework I
% may ask you to work through a short proof of L"ob's theorem, using
% the second incompleteness theorem instead of the fixed-point lemma.

\begin{prob}
نأخذ $\Th{T}$ نظرية مؤكسمة حسابيًا و$\OProv[\Th{T}]$
محمول !!a{derivability} للنظرية $\Th{T}$.
قارن العبارات الأربع الآتية:
\begin{enumerate}
\item من $T \Proves !A$ يلزم $T \Proves \OProv[\Th{T}](\gn{!A})$.
\item $T \Proves !A \lif \OProv[\Th{T}](\gn{!A})$.
\item من $T \Proves \OProv[\Th{T}](\gn{!A})$ يلزم $T \Proves !A$.
\item $T \Proves \OProv[\Th{T}](\gn{!A}) \lif !A$
\end{enumerate}
عيّن شروط صدق كل عبارة منها.
\end{prob}

\end{document}
